% TOP_PROBLEM: {"id": 2603, "title": "Kourovka Notebook Problem 21.94", "category_id": 17, "category": {"id": 17, "name": "group_theory", "display_name": "Group Theory", "description": "Problems about groups, group actions, representations, and related algebraic structures.", "slug": "group-theory", "order_index": 17, "created_at": "2026-07-31T15:26:25.671Z"}, "status": "open", "problem_number": "KOU-21.94", "set_id": 10}
% FIRST_POSTED: 2026-10-01
% ENTRY_KIND: statement_only
% UnsolvedMath ID 2603; source catalogue code KOU-21.94
% !TeX program = lualatex
% Definition and sources only; this file is not an LLM solution attempt.
\documentclass[11pt,a4paper]{article}
\usepackage[margin=25mm]{geometry}
\usepackage{fontspec}
\setmainfont{lmroman10-regular.otf}[BoldFont=lmroman10-bold.otf,
 ItalicFont=lmroman10-italic.otf,BoldItalicFont=lmroman10-bolditalic.otf]
\usepackage{amsmath,amssymb,mathtools}
\usepackage{unicode-math}
\setmathfont{latinmodern-math.otf}
\AtBeginDocument{\let\setminus\smallsetminus}
\usepackage{xurl,hyperref}
\hypersetup{colorlinks=true,urlcolor=blue}
\setcounter{secnumdepth}{0}
\setlength{\parindent}{0pt}
\setlength{\parskip}{5pt}
\setlength{\emergencystretch}{3em}
\begin{document}
\section*{Kourovka Notebook Problem 21.94}\label{2603}
{\small UnsolvedMath ID 2603 · KOU-21.94 · Group Theory · Kourovka Notebook - New Problems, Issue 21}
\subsection{Definitions and mathematical statement}
For a finite group $G$, let $\pi(G)$ be the set of primes dividing $|G|$. Its Gruenberg--Kegel graph, or prime graph, has vertex set $\pi(G)$ and an edge between distinct primes $p,q$ exactly when $G$ contains an element of order $pq$. There are no loops or multiple edges.

Write $\overline{\operatorname{GK}}(G)$ for the abstract graph obtained by forgetting which prime labels each vertex. An isomorphism of such graphs is any adjacency-preserving bijection between their vertex sets; it need not preserve prime labels. Call $G$ recognizable by its unlabelled prime graph if
\[
 \forall\text{ finite groups }H,\qquad
 \overline{\operatorname{GK}}(H)\cong
 \overline{\operatorname{GK}}(G)\ \Longrightarrow\ H\cong G.
\]
The question asks whether infinitely many pairwise nonisomorphic finite groups have this property. Equivalently, are there groups $G_1,G_2,\ldots$ with $G_i\not\cong G_j$ for $i\ne j$, each uniquely determined among all finite groups by the abstract isomorphism type of its prime graph?

The possible competitors $H$ are unrestricted finite groups: they are not assumed to have the same order, the same prime divisors, or the same composition factors as $G$. In particular the recognition condition must exclude a competitor whose prime graph has the same shape but different vertex labels. Group isomorphism, rather than equality of presentations, is the required conclusion.
\subsection{Short English statement}
Are infinitely many finite groups uniquely determined, among all finite groups, by the shape of their unlabelled prime graphs?
\subsection{Sources}
\begin{itemize}
\item[S1] ULAM.AI and the UnsolvedMath Contributors, supplied problems.json, record 2603 (KOU-21.94), Kourovka Notebook - New Problems, Issue 21. Catalogue identity and source transcription; CC BY 4.0.
\url{https://huggingface.co/datasets/ulamai/UnsolvedMath}.
\item[S2] The Kourovka Notebook, 21st edition, version 47 (30 September 2026), new problems, Problem 21.94. Expanded formulation of the supplied catalogue statement.
\url{https://arxiv.org/pdf/1401.0300v47}.
\end{itemize}
\end{document}
